

\bsec{Fixpunkte}{
		
		\begin{itemize}
			\item Newton-Verfahren
			\item Fixpunkte
			\item Banach'scher Fixpunktsatz
			\item Verständnis/Trick Fragen
		\end{itemize}
		}
		
%----------------------------------------------------------------------------------
\bsubsec{Newton-Verfahren}{
		
		Newton-Verfahren:
		$$x_{n+1} = x_n - \frac{f(x_n)}{f'(x_n)}$$
		
	}

	\btask{
		\begin{figure}[ht]
			\label{pic:fixpunkt}
			\centering
			\begin{tikzpicture}
			\datavisualization [school book axes,
			visualize as smooth line,
			y axis={label={$f(x)$}},
			x axis={label} ]
			
			data [format=function] {
				var x : interval [-2:4] samples 22;
				func y = .5 *((\value x*\value x*\value x) / 4 - (3 * \value x));
			};
			\end{tikzpicture}
		\end{figure}
		Mal das Bild vielleicht ab. Du wirst es in den nächsten Aufgaben benötigen und brauchst dann nicht immer hochzuscrollen.
	}

	\btask{

		Alle Aufgaben beziehen sich auf obige Grafik auf Seite \ref{pic:fixpunkt}.
		
		\begin{enumerate}
			\item Nur aus der Grafik gelesen: Wenn man das Newton-Verfahren für diese Funktion anwenden würde, entspräche $x=0$ einem anziehenden Fixpunkt oder nicht? Warum?
			\item Zeige deine Theorie, indem du $x_0 = -1$ als Startwert nimmst und drei Schritte des Newton-Verfahrens grafisch ausführst.
			\item Jemand behauptet, mit Newton wäre die Iterationsvorschrift für die Funktion $f(x)$:
			$$x_{n+1} = x_n^2 - 9$$
			Warum kann das nicht stimmen?
		\end{enumerate}
		
	}
	\bsol{
		
		\begin{enumerate}
			\item Ja, weil wenn man in der Nähe dieser Nullstelle Werte als Startwert von Newton nimmt, würde das Verfahren mit wenigen Schritten die Nullstelle erreicht haben.
			\item Spätestens beim dritten Schritt ist man sehr nahe bei $x=0$
			\item Die Schnittstellen der Iterationsfunktion und $f(x)$ entsprechen den Fixpunkten, die damit ausgerechnet werden sollen. $x=0$ und $x \sim 3.46$ sind weder Fixpunkte dieser Funktion noch Schnittstellen mit der Funktion $f(x)$.
		\end{enumerate}
	}
%----------------------------------------------------------------------------------
\bsubsec{Fixpunkte}{
		
		$x*$ ist Fixpunkt einer Funktion $f(x)$, genau dann wenn $f(x*) = x*$
		
		Iterationsvorschrift $\Phi(x)$
		$$x_{i+1} = \Phi(x_i)$$
%		Anziehender Fixpunkt (konvergiert) $\longleftrightarrow$
%		$$\abs{\Phi'(x*)} < 1$$
%		Abstoßender Fixpunkt (divergiert):
%		$$\abs{\Phi'(x*)} > 1$$
%		Indifferent:
%		$$\abs{\Phi'(x*)} = 1$$
		$\abs{\Phi'(x*)} < 1 \Longrightarrow$ anziehender Fixpunkt (konvergiert) \\
		$\abs{\Phi'(x*)} > 1 \Longrightarrow$ abstoßender Fixpunkt (divergiert) \\
		$\abs{\Phi'(x*)} = 1 \Longrightarrow$ indifferent
	}

	\btask{
		\begin{figure}[ht]
			\label{pic:fixpunkt2}
			\centering
			\begin{tikzpicture}
			\datavisualization [school book axes,
			visualize as smooth line,
			y axis={label={$g(x)$}},
			x axis={label} ]
			
			data [format=function] {
				var x : interval [-0.5:2.5] samples 22;
				func y = (\value{x} - 1.0) * (\value{x} - 1.0) * (\value{x} - 1.0) - \value{x} + 2;
			};
			\end{tikzpicture}
			%((x-1)^3 - x + 2)
		\end{figure}
		Mal das Bild vielleicht ab. Du wirst es in den nächsten Aufgaben benötigen und brauchst dann nicht immer hochzuscrollen.
	}
	\btask{
		
		
		\begin{enumerate}
			\item Zeige, wie du grafisch auf alle Fixpunkte der Funktion $g(x)$ kommen würdest.
			\item Beweise deine Theorie für einen Fixpunkt rechnerisch. Nimm hierfür an, dass 
				$$g(x) = (x-1)^3 - x + 2$$
			\item Eignet sich folgende Iterationsvorschrift theoretisch, um $-\frac{1}{a}$ auszurechnen?
				$$\Phi(x) = \frac{a + a^2 \cdot x + x^2}{x}$$
		\end{enumerate}
		
	}
	\bsol{
		
		\begin{enumerate}
			\item Zeichne die Gerade $m(x) = x$ ein. Alle Schnittpunkte zwischen $m(x)$ und $g(x)$ entspricht den Fixpunkten von $g(x)$. 
			\item $$g(1) = (1-1)^3 - 1 + 2 = -1 + 2 = 1$$
			\item Fixpunkt $-\frac{1}{a}$ einsetzen:
				$$\Phi(-\frac{1}{a}) = \frac{a + a^2 \cdot -\frac{1}{a} + (-\frac{1}{a})^2}{-\frac{1}{a}} = = (-\frac{1}{a})^2 \cdot (-a) = -\frac{1}{a}$$ Passt. Jetzt noch überprüfen, ob dieser anziehend ist:
				$$\Phi'(x) = \frac{x^2 - a}{x^2}$$
				$$|\Phi'(-\frac{1}{a})| = |\frac{(-\frac{1}{a})^2 - a}{(-\frac{1}{a})^2}| = |\frac{(-\frac{1}{a})^2}{(-\frac{1}{a})^2} + \frac{-a}{(-\frac{1}{a})^2}| = |1 - a^3|$$
				Daraus folgt: Genau dann wenn $-1 \le a \le 0$, eignet sich die Vorschrift, ansonsten nicht. (In der Praxis dadurch natürlich Quatsch!)
		\end{enumerate}
	}

%----------------------------------------------------------------------------------
\bsubsec{Banach'scher Fixpunktsatz}{
		
		Es existiert ein anziehender Fixpunkt im Intervall $I$ genau dann, wenn
		\begin{enumerate}
			\item $I$ abgeschlossen
			\item $\forall x \in I: \Phi(x) \in I$
			\item $\Phi$ ist ''Lipschitz-stetig'' im Intervall $I$ mit $0 < L < 1$: \\
			$$\exists L: \abs{\Phi(x) - \Phi(y)} \le L \cdot \abs{x - y}$$
		\end{enumerate}
	}
	\btask{
		
		Sei
		$$f(x) = (x-1)^2 - x + 2$$
		\begin{enumerate}
			\item Berechne alle Fixpunkte von $f(x)$ rechnerisch!
			\item Mithilfe des Intervalls $[0.5; 1.5]$, zeige mithilfe des Banach'schen Fixpunktsatzes, ob sich in dem Intervall ein anziehender Fixpunkt befindet oder nicht.
		\end{enumerate}
		
	}
	\bsol{
		
		\begin{enumerate}
			\item $$f(x) = (x-1)^2 - x + 2 = x \leftrightarrow 0 = (x-1)^2 + 2 - 2x$$
			 $$\leftrightarrow 0 = (x^2 - 2x + 1) + 2 - 2x \leftrightarrow 0 = x^2 -4x + 3$$
			 PQ-Formel:
			 $$x_{1,2} = -\frac{-4}{2} \pm \sqrt{(\frac{4}{2})^2 - 3} = 2 \pm 1$$
			 $$x_1 = 1 \aae x_2 = 3$$
			 %$$\leftrightarrow 0 = (x-1)^2 + 2\cdot (1-x)$$
			 \item \begin{enumerate}
			 	\item Trivial
			 	\item Ist streng monoton fallend in dem Intervall. Größter Wert bei 
			 	$$f(\frac{1}{2}) = (-0.5)^2 - 0.5 + 2 = 1.75$$
			 	Kleinster Wert bei
			 	$$f(\frac{3}{2}) = (0.5)^2 - \frac{3}{2} + 2 = \frac{3}{4}$$
			 	Nicht erfüllt! Kein anziehender Fixpunkt in diesem Intervall!
			 \end{enumerate}
			 %2x -3 ist Ableitung
		\end{enumerate}
	}
%----------------------------------------------------------------------------------
\bquestion{
	}
	\btask{
		
		Beantworte folgende (Trick-)Fragen möglichst genau:
		\begin{enumerate}
			\item TBD
		\end{enumerate}
		
	}
	\bsol{
		
		\begin{enumerate}
			\item TBD
		\end{enumerate}
	}
%----------------------------------------------------------------------------------